\section*{Section \ref{S:equivclasses}}

\begin{list}{\bf{\ref{exer:sec73-1}.}}
\item Use the directed graph to examine all the cases necessary to prove that $\sim$ is reflexive, symmetric, and transitive.  The distinct equivalence classes are: $\left[ a \right] = \left[ b \right] = \left\{ a,b \right\}$; \quad 
$\left[ c \right] = \left\{ c \right\}$; \quad
$\left[ d \right] = \left[ e \right] = \left\{ d,e \right\}$
\end{list}


\begin{list}{\bf{\ref{exer:drawdirgraph}.}}
\item The equivalence class are
\[
[a] = [b] = [d] = \{a, b, d \}, \quad [c] = \{ c \}, \quad [e] = [f] = \{ e, f \}.
\]
Following is the directed graph for this equivalence relation.
\begin{figure}[!h]
\begin{center}
\scalebox{0.6}{\includegraphics{figps-exer2-73.eps}}
\end{center}
\end{figure}
\end{list}


\begin{list}{\bf{\ref{exer:sec73-2}.}}
\item Let $x \in A$.  Since $x$ has the same number of digits as itself, the relation $\mathrel{R}$ is reflexive.  Now let $x, y, z \in A$.  If $x \mathrel{R} y$, then $x$  and  $y$  have the same number of digits.  Hence, $y$ and $x$ have the same number of digits and $y \mathrel{R} x$, and so $\mathrel{R}$ is symmetric.

If $x \mathrel{R} y$ and $y \mathrel{R} z$, then $x$  and  $y$  have the same number of digits and 
$y$  and  $z$  have the same number of digits.  Hence, $x$  and  $z$  have the same number of digits, and so $x \mathrel{R} z$.  Therefore, $\mathrel{R}$ is transitive.

The equivalence classes are:  
$\left\{ 0, 1, 2, \ldots , 9 \right\}$, $\left\{ 10, 11, 12, \ldots , 99 \right\}$, \\
$\left\{ 100, 101, 102, \ldots , 999 \right\}$, $\left\{ 1000 \right\}$.
\end{list}



\begin{list}{\bf{\ref{exer:sec73-congclass}.}}
\item The congruence classes for the relation of congruence modulo 5 on the set of integers are

\begin{multicols}{2}$\left[ 0 \right] = \left\{ 5n \mid n \in \mathbb{Z} \right\}$

$\left[ 1 \right] = \left\{ 5n +1 \mid n \in \mathbb{Z} \right\}$

$\left[ 2 \right] = \left\{ 5n + 2 \mid n \in \mathbb{Z} \right\}$

$\left[ 3 \right] = \left\{ 5n + 3 \mid n \in \mathbb{Z} \right\}$

$\left[ 4 \right] = \left\{ 5n + 4 \mid n \in \mathbb{Z} \right\}$.
\end{multicols}
\end{list}


\begin{list}{\bf{\ref{exer:modpowers}.}}
\item \begin{list}{\bf{(a)}}
\item Let $a, b, c \in \Z_9$.  Since $a^2 \equiv a^2 \pmod 9$, we see that $a \sim a$ and $\sim$ is reflexive.  Let $a, b, c \in \Z_9$.  If $a \sim b$, then $a^2 \equiv b^2 \pmod 9$ and hence, by the symmetric property of congruence, $b^2 \equiv a^2 \pmod 9$.  This proves that $\sim$ is symmetric.  Finally, if $a \sim b$ and $b \sim c$, then 
$a^2 \equiv b^2 \pmod 9$ and $b^2 \equiv c^2 \pmod 9$.  By the transitive property of congruence, we conclude that $a^2 \equiv c^2 \pmod 9$ and hence, $a \sim c$.  This proves that $\sim$ is transitive.  The distinct equivalence classes are $\left\{0, 3, 6 \right\}$, $\left\{1, 8 \right\}$, 
$\left\{ 2, 7 \right\}$, and $\left\{ 4, 5 \right\}$.
\end{list}
\end{list}


\begin{list}{\bf{\ref{exer:sec73-integer}.}}
\item \begin{list}{\bf{(a)}}
\item Let $x \in \left[ \dfrac{5}{7} \right]$.  Then $x - \dfrac{5}{7} \in \Z$, which means that there is an integer $m$ such that $x - \dfrac{5}{7} = m$, or $x = \dfrac{5}{7} + m$.  This proves that \linebreak
$x \in \left\{ \left. m + \dfrac{5}{7} \right| m \in \Z \right\}$ and, hence, that
$\left[ \dfrac{5}{7} \right] \subseteq \left\{ \left. m + \dfrac{5}{7} \right| m \in \Z \right\}$.  We still need to prove that 
$\left\{ \left. m + \dfrac{5}{7} \right| m \in \Z \right\} \subseteq \left[ \dfrac{5}{7} \right]$.
\end{list}
\end{list}


\begin{list}{\bf{\ref{exer:sec73-rationals}.}}
\item \begin{list}{\bf{(a)}}
\item To prove the relation is symmetric, note that if $(a, b) \approx (c, d)$, then 
$ad = bc$.  This implies that $cb = da$ and, hence,  $(c, d) \approx (a, b)$.
\end{list}
\end{list}


\begin{list}{}
\item \begin{list}{\bf{(c)}}
\item $3a = 2b$
\end{list}
\end{list}

\hbreak
\endinput


